\documentclass[10pt]{article}
\usepackage[margin=0.76in]{geometry}
\usepackage{amsmath,amssymb,amsthm,bm,booktabs,graphicx,microtype}
\usepackage{hyperref,xcolor}
\newtheorem{theorem}{Theorem}
\newtheorem{remark}{Remark}
\DeclareMathOperator{\softmax}{softmax}
\title{Scaling Laws for In-Context Bayesian Near-Field Linear Sampling}
\author{Technical results note}
\date{5 August 2026}
\begin{document}
\maketitle

\begin{abstract}
We study how an in-context posterior-moment solver for the original two-
dimensional near-field linear sampling method scales with the number $n$ of
training inverse problems, network size $P$, context length $m$, recurrent
depth, and wall-clock time.  The input is a complex multistatic response from
deterministic point sources around one to six sound-soft obstacles.  A fixed
von Mises/softmax covariance selects the Gaussian-process feature space; its
nonlinearity is a modelling choice and is not learned.  Richardson, heavy-ball
(HB), Chebyshev and PCG share the same context encoder and Bayesian decoder and
solve posterior mean and covariance right-hand sides in parallel.  Identity-
preconditioned CG is the matched-matrix-vector-product control.  The sweep
contains 2 trained checkpoints and 336 learned-
model/task evaluations.  We report task-level uncertainty, inference time,
empirical scaling exponents and a simultaneous distribution-free held-out
risk bound.  The main question is not whether a large controller can imitate
CG, but when learned conditioning improves CG under severe, task-dependent
ill-conditioning.
\end{abstract}

\section{Problem and posterior moments}
Let $N_D\in\mathbb C^{m\times m}$ be the whitened near-field matrix for a
sound-soft obstacle $D$, $\Phi=[\phi_z]_z$ the Green-function probes on a
two-dimensional sampling grid, and $K_\Gamma\succ0$ a prescribed angular GP
covariance.  The nonlinear feature choice is
\[
 W_{ij}=\exp\{\gamma\cos(\theta_i-\theta_j)\},\qquad
 K_\Gamma=(1-\eta)I+\eta D_W^{-1/2}WD_W^{-1/2},
\]
with fixed $\gamma=1$ and $\eta=0.2$.  Define
\[
 H_D=N_DK_\Gamma N_D^*+I,\qquad
 H_DQ_\mu=\Phi,\qquad H_DQ_\Sigma=N_DK_\Gamma.
\]
One tied recurrence acts on $[\Phi,N_DK_\Gamma]$ and yields
\[
 M_D=K_\Gamma N_D^*Q_\mu,\qquad
 \Sigma_D=K_\Gamma-K_\Gamma N_D^*Q_\Sigma.
\]
Thus a small mean residual is insufficient for Bayesian correctness: the
covariance residual must converge as well.

\section{Variable-context architecture and matched solvers}
The covariance block uses a deterministic orthonormal harmonic sketch with a
fixed number of columns and $m$ rows.  Its feature dimension is therefore
independent of context length, while the physical operator is never padded or
replaced by a learned surrogate.  Training contexts are
[12, 16, 24, 32]; evaluation contexts are
[8, 12, 16, 24, 32, 48], so the endpoints include context-length
extrapolation.  Controller widths are [128] and sample
sizes are [128, 256].

At test depth $T=32$, all recurrent methods use exactly
$T$ applications of the transformed operator for both posterior blocks.
Identity-CG sets the population factor to $I$.  Population-PCG learns a positive
geometry/population factor and retains exact task-wise Krylov coefficients.
Looped HB, Chebyshev and Richardson additionally infer finite-horizon spectral
endpoints from RHS-weighted Krylov statistics.  Consequently, PCG versus CG
isolates learned conditioning; a looped method versus PCG measures whether a
stationary learned recurrence is worth replacing Krylov adaptation.

\section{Axes and scaling model}
The dataset size $n$ is the number of unique physical inverse problems
presented once in an online stream; the eight examples in an optimizer batch
share one sampled acquisition configuration.  The context length $m$ is both
the number of point sources and receivers, giving $m^2$ complex measurements.
Network size $P$ counts trainable parameters.  Each held-out configuration
contains fixed, never-trained MFS tasks spanning ID obstacles and OOD obstacle
count, shape, noise, aperture and frequency.

For localization risk $R_{\rm loc}=1-\mathrm{AP}$ and bounded numerical
risk $R_{\rm solve}=r/(1+r)$, we separately fit the descriptive law
\begin{equation}
 R(n,m,P)=R_\infty+A(n/n_0)^{-\alpha}+B(m/m_0)^{-\beta}
 +C(P/P_0)^{-\gamma}+D\sqrt{(P/P_0)/(n/n_0)}.
 \label{eq:scaling}
\end{equation}
The last term allows a finite-data width penalty.  Equation~\eqref{eq:scaling}
is an empirical scaling model, not a theorem.

\begin{table}[t]
\centering\small
\caption{Empirical exponents in Eq.~\eqref{eq:scaling} for both risks;
larger positive values indicate faster improvement along that axis.}
\begin{tabular}{llcccc}
\toprule Method & risk & data $\alpha$ & context $\beta$ & width $\gamma$ & $R^2$\\
\midrule
population-PCG & localization & 0.010 & 3.000 & 2.194 & 0.984 \\
population-PCG & solver & 0.010 & 0.034 & 0.213 & -0.096 \\
\bottomrule
\end{tabular}
\end{table}

\section{A simultaneous held-out generalization bound}
Let $L_c\in[0,1]$ be localization risk for a pre-specified trained
configuration $c=(n,m,P,$ method$)$, and let $\widehat R_c$ average $T_c$
independent held-out tasks.  A union bound and one-sided Hoeffding inequality
give the following non-asymptotic certificate.
\begin{theorem}[Simultaneous task-distribution bound]
For $K$ evaluated configurations, with probability at least $1-\delta$,
simultaneously for every $c$,
\[
 R_c\leq \widehat R_c+
 \sqrt{\frac{\log(K/\delta)}{2T_c}}.
\]
\end{theorem}
We use $\delta=0.05$ and report the clipped upper bound in
Figure~\ref{fig:bounds}.  Training samples determine the learned hypothesis
and hence $\widehat R_c$; held-out sample size controls certification slack.
The same theorem is applied to the bounded numerical risk
$L_c^{\rm solve}=r_c/(1+r_c)$, where $r_c$ is the posterior-mean relative
residual.  This second risk remains sensitive after localization AP saturates.
This distinction matters.  Replacing $T_c$ by training size $n$ after fitting
would be invalid, while a generic parameter-count uniform bound for this
nonconvex complex network is vacuous at the present $P/n$ ratios.

\begin{figure}[t]
\centering
\includegraphics[width=0.98\linewidth]{scaling_generalization_bounds.png}
\caption{Empirical localization and numerical solver risks with simultaneous
95\% held-out upper bounds across dataset and context sizes.}
\label{fig:bounds}
\end{figure}

\section{Results}
\begin{table}[t]
\centering\small
\caption{Final comparison at context $m=24$, width
128 and $n=256$.  Baselines do
not train.  Runtime is milliseconds per held-out batch.}
\begin{tabular}{lccccc}
\toprule Method & AP & mean res. & covariance res. & score error & ms\\
\midrule
population-PCG & 0.768 & 0.213 & 6.25e-03 & 3.97e-03 & 21.291 \\
identity-CG & 0.769 & 0.213 & 6.50e-03 & 4.11e-03 & 15.106 \\
exact & 0.768 & 5.98e-05 & 1.65e-06 & 0 & 10.968 \\
\bottomrule
\end{tabular}
\end{table}

\begin{figure}[t]
\centering
\includegraphics[width=0.98\linewidth]{scaling_dataset.png}
\caption{Learning curves in dataset size at fixed network width and context.}
\end{figure}

\begin{figure}[t]
\centering
\includegraphics[width=0.98\linewidth]{scaling_context.png}
\caption{Context-length scaling.  Context $m=8$ and $m=48$ are outside the
training context set.}
\end{figure}

\begin{figure}[t]
\centering
\includegraphics[width=0.98\linewidth]{scaling_network.png}
\caption{Network-size scaling at the largest completed dataset size.}
\end{figure}

\begin{figure}[t]
\centering
\includegraphics[width=0.9\linewidth]{scaling_time.png}
\caption{Wall-clock learning curves.  Time includes physical task generation,
forward/backward propagation and optimization, but excludes held-out audits.}
\end{figure}

\begin{figure}[t]
\centering
\includegraphics[width=0.95\linewidth]{scaling_runtime.png}
\caption{Inference scaling with $m^2$ complex measurements and the residual--
time Pareto comparison at $m=24$.}
\end{figure}

The equal-budget residual ratios at the final central configuration are
1.004$\times$ for
the posterior mean and
1.039$\times$
for posterior covariance.  Ratios above one isolate a benefit from the learned
population factor beyond ordinary CG.  Localization can nevertheless saturate
before residual convergence; in that regime additional solver accuracy does
not overcome the physical point-spread limit.

\section{Interpretation and limitations}
Three claims are deliberately separated.  First, CG is a strong task-adaptive
baseline because its coefficients already depend on the current Krylov space.
Second, population-PCG tests whether geometry-distribution conditioning adds
value when no obstacle-specific preconditioner is known.  Third, looped HB
tests amortized fixed-depth iteration, not a universally superior replacement
for CG.  A true task-specific ICL preconditioner would additionally map current
moment or Ritz information into an SPD operator
$C_D(A_D)\approx A_D^{-1/2}$; the present population factor does not rotate
all obstacle-dependent eigendirections.

The held-out bound certifies the stated simulator distribution only.  The
forward data are generated by MFS, not laboratory measurements, and context
extrapolation does not imply robustness to arbitrary receiver geometries.
Physical resolution remains controlled by aperture, frequency and sensor
count.  These qualifications prevent a small numerical loss from being
misreported as solved inverse scattering.

\section{Reproducibility}
All task-level records, cumulative training times, checkpoints, runtime
measurements, scenario definitions, fitted parameters and bound values are
stored beside this PDF.  The experiment uses three fixed seeds and preserves
the exact same held-out tasks across methods, widths and dataset checkpoints.

\begin{thebibliography}{9}
\bibitem{kang2026} G. Kang, C. J. Lee, and X. Cheng.
Transformers can learn posterior predictive distributions in-context.
\emph{arXiv:2605.26713}, 2026.
\bibitem{garnier2022} J. Garnier, H. Haddar, and H. Montanelli.
The linear sampling method for random sources. \emph{arXiv:2210.15560}, 2022.
\bibitem{garnier2024} J. Garnier, H. Haddar, and H. Montanelli.
The linear sampling method for data generated by small random scatterers.
\emph{arXiv:2403.19482}, 2024.
\bibitem{polyak1964} B. T. Polyak. Some methods of speeding up the convergence
of iteration methods. \emph{USSR Computational Mathematics and Mathematical
Physics}, 1964.
\end{thebibliography}
\end{document}
